\newpage
\thispagestyle{empty}

\begin{center}
{\Large\textbf{ABSTRACT}}
\end{center}

\vspace{1cm}

The administrative process of managing campus hostel facilities involves tracking 
room allocations, student records, fee payments, and maintenance requests. 
Maintaining this information using scattered physical ledgers, disconnected 
spreadsheets, or other manual methods makes it difficult to track real-time 
occupancy and operational status. To address this problem, this project proposes 
the \textbf{Hostel Administration System}, a web-based management platform 
designed to provide administrators and managers with a centralised digital 
environment for organising and monitoring hostel accommodations.

The system allows authorised users to securely log in and manage hostel-related 
information through an authenticated interface. It provides robust functionality 
for recording and updating details such as hostel blocks, room capacities, 
current occupancies, and specific student allocations. By maintaining a 
structured approach to data, the system ensures that room availability and 
student records are tracked accurately and consistently, replacing manual 
cross-referencing with automated tracking logic.

The system is implemented using a modern full-stack web application architecture 
(MERN stack). The frontend is developed using \textbf{React} with \textbf{Vite}, 
while \textbf{Tailwind CSS} provides a responsive user interface. \textbf{React 
Router} is utilised for client-side navigation and protecting administrative routes. 
The backend is powered by a \textbf{Node.js} and \textbf{Express.js} API. User 
authentication is implemented using \textbf{JSON Web Tokens (JWT)}. Administrative 
records and student data are persistently stored using a \textbf{MongoDB} database.

The system follows a modular design in which the presentation layer, routing, 
backend controllers, and database schemas are maintained independently. This 
approach improves maintainability and establishes a foundation for extending the 
application with features such as automated fee tracking, a dedicated student 
complaint portal, and occupancy statistics visualisation. 

The Hostel Administration System therefore provides a structured, secure approach 
to managing campus accommodations, significantly reducing the difficulty and error 
rates associated with manual tracking. By combining an intuitive user interface 
with robust backend API services and scalable cloud storage, the system provides 
a centralised and accessible platform for managing hostel operations.
